\section{Quarkus – Grundlagen und Einsatz in der Anwendung}
\setauthor{Markus Tran}
Quarkus ist ein modernes, für den Einsatz in containerisierten und Cloud-Umgebungen optimiert JavaFramework.
Es wurde entwickelt, um den Prozess des Deployment von Java-Applikationen in Cloud-Umgebungen zu verbessern.
Das Framework unterstützt Entwickler dabei, auch ohne tiefgehende Kenntnisse der zugrundeliegenden Infrastruktur, ihre Applikationen der Cloud-Umgebung bereitzustellen.
Quarkus ist darauf basiert, mit bereits existierenden Frameworks, wie 
z.B. "Eclipse MicroProfile und Spring sowie für Apache Kafka, RESTEasy (JAX-RS), Hibernate ORM (JPA), Spring, Infinispan, Camel "
gut zusammen zu arbeiten.
Durch die Leichtgewichtigkeit hebt sich Quarkus von anderen Frameworks ab. Dies wird realisiert durch Faktoren wie, durch eine geringe Startzeit,
geringen Ressourcenverbrauch. Doch ist der größte Aspekt der für Quarkus spricht, die Möglichkeit, Projekte in Cloud-Umgebungen bereitzustellen auch ohne großen Aufwand. 

\section{Grundprinzip von Quarkus}
\setauthor{Markus Tran}
Was Quarkus hervorhebt, ist die Art wie es Java-Anwendungen kompiliert und ausführt.
Java unterstützt nativ die Just-In-Time Kompilierung (JIT-Kompilierung), welche Codeblöcke in Assembler-Code übersetzt. 
Das heißt es wandelt den vom Entwickler geschriebenen Code, der für Menschen lesbar ist, in Assembler-Code. 
Denn, erst der Assembler Code ist für die Maschiene lesbar.
Im Gegensatz dazu verwendet Quarkus Ahead-of-Time (AOT) Kompilierung, um den gesammten Code vor der Ausführung in eine native ausführbare Datei umzuwandeln. 

\subsection{JIT-Kompilierung}

\subsection{Ahead-of-Time}

\subsection{Hot Reload}

\subseaction{Kubernetes}

\subseaction{Dependency Injection}

\section{Quarkus in der vorliegenden Applikation}
\setauthor{Markus Tran}
Quarkus wird in der vorliegenden Anwendung als Hauptframework für das Backend verwendet. 
Die vorliegende Applikation verwendet Hibernate ORM (Panache) für die Kommunikation und das Verwalten der Datenbank.
Wie Ahead-Of-Time(AOT) Kompilierung in diesem Falle nützlich ist trotz der längerer Build-Zeit im Vergleich zu Just-In-Time(JIT) Kompilierung. 
Druch das kompilieren aller Datein, während des Build Prozesseses, muss nicht während der Laufzeit der Applikation, weiter optimiert werden.
Die Anwendung weist dadurch, einen geringen Speicherbedarf, schnelleren Start von Services und eine containerisierte Umgebung, auf.


\subsection{Deeper}
Nicht mehr im Inhaltsverzeichnis.

\subsubsection{Deepest}
Vermeide mich.